\chapter{Conception de la solution}

\section{Architecture logique}

L'architecture adopte un modèle \textbf{client--serveur} à trois composants principaux :
\begin{itemize}[leftmargin=*]
  \item une \textbf{API REST} monolithique (\textbf{Node.js} / \textbf{Express}) ;
  \item une application \textbf{Flutter} (utilisateurs) ;
  \item une interface \textbf{React} / \textbf{Vite} (administration).
\end{itemize}

La logique serveur suit une organisation \textbf{Routes $\rightarrow$ Contrôleurs $\rightarrow$ Services}, avec des \textbf{middlewares} (authentification JWT, gestion d'erreurs, limitation de débit, téléversement). Les données sont stockées dans \textbf{MySQL} (moteur InnoDB, encodage \texttt{utf8mb4}).

\section{Diagramme de composants}

La figure~\ref{fig:components} présente la répartition en couches et les dépendances externes (Firebase, OAuth Google, fournisseur de cartographie, CI GitHub Actions).

\begin{figure}[H]
\centering
\IfFileExists{figures/component_architecture_moroccocheck.png}{%
  \includegraphics[width=\linewidth]{figures/component_architecture_moroccocheck.png}%
}{%
  \placeholderfig{Diagramme de composants — architecture}{figures/component\_architecture\_moroccocheck.png}%
}
\caption{Diagramme de composants — architecture MoroccoCheck}
\label{fig:components}
\end{figure}

\section{Modèle de données}

Les entités centrales incluent notamment : \textbf{User}, \textbf{TouristSite}, \textbf{Category}, \textbf{CheckIn}, \textbf{Review}, \textbf{Badge}, \textbf{UserBadge}, \textbf{ContributorRequest}, ainsi que des tables de liaison et de session. Les \textbf{vues SQL} (ex.\ \texttt{v\_user\_stats}, \texttt{v\_site\_details}) facilitent les statistiques ; des \textbf{déclencheurs} maintiennent la cohérence des compteurs et des moyennes de notes lors des insertions ou suppressions d'avis et de check-ins.

\section{Diagramme de classes}

Le diagramme de classes (figure~\ref{fig:classes}) synthétise les associations principales entre entités du domaine.

\begin{figure}[H]
\centering
\IfFileExists{figures/class_diagram_moroccocheck.png}{%
  \includegraphics[width=\linewidth]{figures/class_diagram_moroccocheck.png}%
}{%
  \placeholderfig{Diagramme de classes du domaine}{figures/class\_diagram\_moroccocheck.png}%
}
\caption{Diagramme de classes — domaine métier MoroccoCheck}
\label{fig:classes}
\end{figure}

\section{Cas d'utilisation}

La figure~\ref{fig:uc} illustre les cas d'utilisation globaux et la généralisation entre touriste et contributeur.

\begin{figure}[H]
\centering
\IfFileExists{figures/use_case_global_moroccocheck.png}{%
  \includegraphics[width=\linewidth]{figures/use_case_global_moroccocheck.png}%
}{%
  \placeholderfig{Diagramme de cas d'utilisation global}{figures/use\_case\_global\_moroccocheck.png}%
}
\caption{Diagramme de cas d'utilisation global}
\label{fig:uc}
\end{figure}

\section{Scénarios dynamiques}

\subsection{Authentification}

La figure~\ref{fig:seq-auth} modélise le scénario de connexion : validation des entrées, recherche utilisateur, comparaison de mot de passe (\texttt{bcrypt}), émission de jetons JWT et persistance de session.

\begin{figure}[H]
\centering
\IfFileExists{figures/sequence_auth_moroccocheck.png}{%
  \includegraphics[width=0.95\linewidth]{figures/sequence_auth_moroccocheck.png}%
}{%
  \placeholderfig{Séquence authentification}{figures/sequence\_auth\_moroccocheck.png}%
}
\caption{Diagramme de séquence — authentification (login)}
\label{fig:seq-auth}
\end{figure}

\subsection{Check-in GPS}

La figure~\ref{fig:seq-checkin} décrit le scénario de validation d'un check-in : acquisition GPS côté mobile, appel API authentifié, calcul de distance, règles métier, persistance et mise à jour de la gamification.

\begin{figure}[H]
\centering
\IfFileExists{figures/sequence_checkin_gps_moroccocheck.png}{%
  \includegraphics[width=0.95\linewidth]{figures/sequence_checkin_gps_moroccocheck.png}%
}{%
  \placeholderfig{Séquence check-in GPS}{figures/sequence\_checkin\_gps\_moroccocheck.png}%
}
\caption{Diagramme de séquence — check-in GPS}
\label{fig:seq-checkin}
\end{figure}

\section{Activité — gamification}

Le traitement des points, du rang et des badges est schématisé figure~\ref{fig:act-gamif}. Les seuils de rang dans l'implémentation sont paramétrés dans le code (\texttt{RANK\_THRESHOLDS}) ; le rapport académique peut les aligner sur cette source lors de la rédaction finale.

\begin{figure}[H]
\centering
\IfFileExists{figures/activity_gamification_moroccocheck.png}{%
  \includegraphics[width=\linewidth]{figures/activity_gamification_moroccocheck.png}%
}{%
  \placeholderfig{Diagramme d'activité — gamification}{figures/activity\_gamification\_moroccocheck.png}%
}
\caption{Diagramme d'activité — gamification (points, rangs, badges)}
\label{fig:act-gamif}
\end{figure}
